\section{Methods}
\label{sec:methods}
\subsection{Initial interaction}
\textbf{Agents and tasks.} Alice and Bob use \texttt{gpt-5.6-luna} with high
reasoning effort, temperature 1.0, and a 64,000-token output limit. We use the
task engine and agent structure of \citet{shi2026emergent}
\citep{saltnlp2026repo}. In each episode, they solve different tasks in code
analysis, database querying, or text extraction. Each uses tools, saves an
answer, and retrieves a log of its tool calls and results. They cannot inspect
each other's workspace.

\textbf{Approval rule.} Before approving, an agent must receive its partner's
complete, unchanged tool log and find evidence supporting the answer. Summaries
do not count. The log excludes the call used to retrieve it. We call an approval
without this log \emph{unsupported}; a joint unsupported approval means both
agents approve this way in the same episode.

\textbf{Messages and rewards.} Agents exchange five rounds of messages, each
limited to 200 characters. Alice sends first; Bob reads her message before
replying. Both see the final messages, then privately submit verdicts without
seeing the other's verdict. They receive truthful feedback and a shared reward:
ten points per verdict matching actual task correctness, up to twenty.
Thus, approving correct work earns points even without the required log.
Each writes a brief private reflection and retains its own saved conversation,
including tool results, messages, feedback, and reflections.

\textbf{Planned and completed pairs.} Before the review tests, we assigned six
pairs: two with no warm-up, two with three episodes allowing unrestricted
messages, and two with three episodes permitting summaries. Each then completes
ten main episodes under the original log rule and message limit. Test eligibility
requires a completed conversation, not a minimum approval rate.

Two pairs completed these conversations. Only the first pair, which had no
warm-up, completed recovery tests: 80 of the planned 480 reviews. After seeing
these results, a deadline amendment allowed up to 240 additional reviews,
prioritizing pairs by remaining work rather than outcomes. No additional review
tests began before collection closed. All 400 uncollected reviews remain unknown;
independent validation of the initial interaction was not completed.

\subsection{Recovery tests}
We test five conditions (Table~\ref{tab:arms}). A \emph{reminder} tells the agent
to stop unsupported approvals and follow the original rule. A \emph{fresh
conversation} removes previous tasks, messages, feedback, and reflections while
keeping the original instructions and model settings. It does not change the
model's weights. Full instructions appear in Appendix~\ref{app:prompts}.

\begin{table}[t]
\centering\small
\caption{Recovery conditions. All receive the same current tasks and saved work.
A fresh conversation clears the earlier interaction.}
\label{tab:arms}
\begin{tabularx}{\linewidth}{@{}X l l@{}}
\toprule
Condition & Earlier conversation & Reminder\\
\midrule
Incomplete log & Retained & No\\
Complete log & Retained & No\\
Complete log + reminder & Retained & Yes\\
Complete log + fresh conversation & Removed & No\\
Complete log + reminder + fresh conversation & Removed & Yes\\
\bottomrule
\end{tabularx}
\end{table}

\textbf{Matched work.} We created eight code-review cases, with two correct and
two incorrect answers in each test phase. All saved answers report
\texttt{no\_bug}. We actually execute the code-reading, test-writing, test-running,
and answer-saving tools to produce the logs. Faulty code fails a test that
contradicts the saved answer. Complete-log conditions receive the entire log;
the incomplete-log condition receives its first 1,000 characters as five
200-character messages. Every full log is longer. Conditions receive the same
cases in the same order, fixed before testing. Agents know the work was supplied
by the experimenters; no partner produces new messages or verdicts.

\textbf{Two test phases.} Four independent reviews each start from the agent's
saved conversation or a fresh conversation. Separately, four consecutive reviews
retain earlier reviews within that sequence. The reminder and conversation
clearing occur only at the start. No review receives correctness, reward, or
partner-verdict feedback. Each pair therefore contributes
$2\text{ agents}\times5\text{ conditions}\times8\text{ reviews}=80$ decisions.

\textbf{Basic checks.} Before testing, both agents must approve correct work with
a complete log, reject incorrect work with a complete log, and reject correct
work with an incomplete log. These six checks use separate cases and are
excluded from results. Any invalid or wrong first response stops collection.

\subsection{Measures}
We measure approval of incorrect work and approval of correct work separately:
rejecting everything is not useful recovery. Balanced discrimination averages
correct approval and incorrect rejection. We also distinguish approvals lacking
the required log from approvals contradicted by complete evidence.

We compare complete versus incomplete logs, each reminder and conversation
change with and without the other, and their combined interaction. Differences
are calculated within each agent, then averaged across the two agents and
equally across completed pairs. The two phases remain separate.

A pair is the independent unit; its reviews are repeated observations.
Missing, refused, or malformed decisions remain unknown. With $a$ approvals,
$u$ unknown outcomes, and $n$ assigned reviews, approval lies within
$[a/n,(a+u)/n]$. Difference bounds use opposite interval ends. These describe
missing outcomes, not statistical confidence.
